%% ma: L12-B14-0039
%% lop: 12
%% chuong: 5
%% bai: 14
%% dang: 12.14.19
%% loai: TN4
%% muc: TH
%% dap_an: "B"
%% nguon: "CĐ2.5.pdf giáo viên gửi 27/10/2026 (M07, Lê Minh Khôi – Nguyễn Anh Khôi, Chương II Chủ đề 5 Tích có hướng của hai vectơ), Câu 23"
%% kien_thuc: "Bài 14 (tích có hướng, phương trình mặt phẳng, khoảng cách từ điểm đến mặt phẳng), lớp 11 (thể tích khối chóp, khối hộp)"
%% trang_thai: cho-duyet
%% ly_do: "Dạng mới 12.14.19 đề xuất, chờ giáo viên duyệt"
\begin{ex}
Trong không gian $Oxyz$, cho hình chóp $A.BCD$ có tọa độ các đỉnh $A(0;1;-1)$, $B(1;1;2)$, $C(1;-1;0)$ và $D(0;0;1)$. Độ dài đường cao của hình chóp $A.BCD$ bằng
\choice{$2\sqrt2$}{\True $\dfrac{3\sqrt2}2$}{$3\sqrt2$}{$\dfrac{\sqrt2}2$}
\loigiai{$\left[\overrightarrow{BC},\overrightarrow{BD}\right]=(0;2;-2)$ nên $(BCD)$: $y-z+1=0$.

$d\left(A;(BCD)\right)=\dfrac{|1+1+1|}{\sqrt2}=\dfrac{3\sqrt2}2$.}
\end{ex}
